\documentclass[tikz,border=8pt]{standalone}
\usepackage{ctex}
\usepackage{amsmath}
\usepackage{pgfplots}
\pgfplotsset{compat=1.18}
\usetikzlibrary{calc,positioning,arrows.meta}

%% 图 opt-p2-04-1：一阶 vs 二阶最优性条件三种临界点对比
%% 三个子图：左=极小点（f'=0, f''>0），中=鞍点（f'=0, f''=0），右=极大点（f'=0, f''<0）

\begin{document}
\begin{tikzpicture}

%% ===== 子图 1：局部极小点（碗形）=====
\begin{scope}[xshift=0cm]

  % 坐标轴
  \draw[->, gray!60, thin] (-2.3, 0) -- (2.3, 0)
    node[right, font=\scriptsize] {$x$};
  \draw[->, gray!60, thin] (0, -0.3) -- (0, 3.2)
    node[above, font=\scriptsize] {$f$};

  % 抛物线 f(x) = x^2 + 0.5（碗形）
  \draw[blue!80!black, thick, domain=-1.8:1.8, samples=80]
    plot (\x, {\x*\x + 0.3});

  % 极小点标记
  \fill[red!80!black] (0, 0.3) circle (2.5pt)
    node[below right, font=\scriptsize] {$x^*$};

  % 零梯度箭头（水平切线）
  \draw[red!70!black, thick] (-0.7, 0.3) -- (0.7, 0.3);
  \node[font=\scriptsize, red!70!black] at (1.1, 0.3) {$f'=0$};

  % 正曲率标注（向上弯）
  \draw[<->, green!60!black, thick] (-0.8, 1.0) arc[start angle=200, end angle=340, radius=0.6];
  \node[font=\scriptsize, green!60!black] at (0, 1.6) {$f''>0$};

  % 标题
  \node[font=\small\bfseries, blue!80!black] at (0, -0.8) {极小点};
  \node[font=\scriptsize] at (0, -1.2) {$\nabla f=0,\;\mathbf{H}\succ 0$};

\end{scope}

%% ===== 子图 2：鞍点 / 拐点（平台）=====
\begin{scope}[xshift=5.5cm]

  % 坐标轴
  \draw[->, gray!60, thin] (-2.3, 0) -- (2.3, 0)
    node[right, font=\scriptsize] {$x$};
  \draw[->, gray!60, thin] (0, -0.3) -- (0, 3.2)
    node[above, font=\scriptsize] {$f$};

  % 三次函数 f(x) = x^3 + 1.2（一维鞍点/拐点）
  \draw[orange!80!black, thick, domain=-1.4:1.4, samples=80]
    plot (\x, {\x*\x*\x + 1.2});

  % 驻点（原点处）
  \fill[red!80!black] (0, 1.2) circle (2.5pt)
    node[right, font=\scriptsize] {$x^*$};

  % 零梯度（水平切线）
  \draw[red!70!black, thick] (-0.7, 1.2) -- (0.7, 1.2);
  \node[font=\scriptsize, red!70!black] at (1.15, 1.2) {$f'=0$};

  % 零曲率标注
  \draw[orange!60!black, thick, dashed, domain=-0.9:0.9]
    plot (\x, {1.2});
  \node[font=\scriptsize, orange!60!black] at (0, 2.0) {$f''=0$};
  \node[font=\scriptsize, orange!60!black] at (0, 2.35) {（退化）};

  % 标题
  \node[font=\small\bfseries, orange!80!black] at (0, -0.8) {鞍点（拐点）};
  \node[font=\scriptsize] at (0, -1.2) {$\nabla f=0,\;\mathbf{H}=0$（不定）};

\end{scope}

%% ===== 子图 3：局部极大点（倒碗形）=====
\begin{scope}[xshift=11cm]

  % 坐标轴
  \draw[->, gray!60, thin] (-2.3, 0) -- (2.3, 0)
    node[right, font=\scriptsize] {$x$};
  \draw[->, gray!60, thin] (0, -0.3) -- (0, 3.2)
    node[above, font=\scriptsize] {$f$};

  % 倒抛物线 f(x) = -x^2 + 2.8（倒碗形）
  \draw[purple!70!black, thick, domain=-1.7:1.7, samples=80]
    plot (\x, {-\x*\x + 2.8});

  % 极大点标记
  \fill[red!80!black] (0, 2.8) circle (2.5pt)
    node[above right, font=\scriptsize] {$x^*$};

  % 零梯度（水平切线）
  \draw[red!70!black, thick] (-0.7, 2.8) -- (0.7, 2.8);
  \node[font=\scriptsize, red!70!black] at (1.15, 2.8) {$f'=0$};

  % 负曲率标注（向下弯）
  \draw[<->, purple!60!black, thick] (-0.8, 2.2) arc[start angle=160, end angle=20, radius=0.6];
  \node[font=\scriptsize, purple!60!black] at (0, 1.5) {$f''<0$};

  % 标题
  \node[font=\small\bfseries, purple!70!black] at (0, -0.8) {极大点};
  \node[font=\scriptsize] at (0, -1.2) {$\nabla f=0,\;\mathbf{H}\prec 0$};

\end{scope}

%% ===== 整体标题 =====
\node[font=\large\bfseries] at (7.0, 3.9)
  {一阶 vs 二阶最优性条件：三种临界点对比};

%% ===== 对比说明框 =====
\node[draw=gray!50, thin, fill=gray!5, rounded corners=3pt,
      font=\scriptsize, inner sep=6pt, align=left]
  at (7.0, -1.9)
  {$f'(x^*)=0$（一阶必要条件）对三种临界点均成立\\
   $f''(x^*)>0$ $\Rightarrow$ 极小（二阶充分）\\
   $f''(x^*)<0$ $\Rightarrow$ 极大（二阶充分）\\
   $f''(x^*)=0$ $\Rightarrow$ 不确定（需高阶分析）};

\end{tikzpicture}
\end{document}
